% ============================================================
% Chapter 2 — Linux Preparation Guide
% ============================================================
\chapter{Linux Preparation Guide}
\label{ch:linux}

\noteblock{This chapter targets Kali Linux users. Kali is a Debian-based distribution purpose-built for penetration testing. All commands in this guide assume Kali Linux 2024.x or later.}

% ─────────────────────────────────────────────────────────
\section{Installing Kali Linux}

\subsection{Download}

Download the official Kali Linux ISO from \url{https://www.kali.org/get-kali/}

\begin{table}[H]
\centering
\caption{Recommended Kali Linux installation options}
\begin{tabularx}{\textwidth}{>{\bfseries}lX}
\toprule
Option & Use Case \\
\midrule
Installer ISO & Dedicated machine or dual-boot \\
Virtual Machine (OVA) & VirtualBox or VMware — recommended for beginners \\
WSL2 (Windows) & Lightweight terminal-only option for Windows users \\
Live Boot (USB) & No installation required, runs from USB \\
\bottomrule
\end{tabularx}
\end{table}

\warnblock{Always download from \textbf{kali.org} only. Verify the SHA256 checksum of the ISO before installing. Never trust third-party mirrors for security tools.}

\subsection{Recommended Minimum Specifications}

\begin{table}[H]
\centering
\caption{Minimum hardware for Kali Linux}
\begin{tabular}{ll}
\toprule
Component & Minimum \\
\midrule
RAM & 4 GB (8 GB recommended) \\
Disk & 20 GB (50 GB recommended) \\
CPU & 2 cores (4 cores recommended) \\
Network & Internet access required \\
\bottomrule
\end{tabular}
\end{table}

% ─────────────────────────────────────────────────────────
\section{Updating Kali Linux}

After installation, always update before any testing. Running outdated tools may produce inaccurate results or expose you to vulnerabilities in the tools themselves.

\subsection{Command: \texttt{sudo apt update}}

\begin{lstlisting}[style=terminal, caption={Update the package list}]
sudo apt update
\end{lstlisting}

\begin{table}[H]
\centering
\caption{Command explanation: sudo apt update}
\begin{tabularx}{\textwidth}{>{\bfseries}lX}
\toprule
Field & Detail \\
\midrule
Purpose & Downloads the latest package index from Kali repositories \\
What it does & Refreshes the list of available packages and their versions — does NOT install anything \\
Expected result & Lines like: \texttt{Hit:1 http://kali.download/kali kali-rolling InRelease} \\
Run as & Root (sudo required) \\
\bottomrule
\end{tabularx}
\end{table}

\textbf{Troubleshooting:}
\begin{itemize}
  \item \texttt{Could not resolve 'kali.download'} — check internet connection: \cmd{ping 8.8.8.8}
  \item \texttt{Permission denied} — ensure you prefix with \cmd{sudo}
  \item \texttt{Waiting for lock} — another apt process is running; wait and retry
\end{itemize}

\subsection{Command: \texttt{sudo apt upgrade}}

\begin{lstlisting}[style=terminal, caption={Upgrade installed packages}]
sudo apt upgrade -y
\end{lstlisting}

\begin{table}[H]
\centering
\caption{Command explanation: sudo apt upgrade}
\begin{tabularx}{\textwidth}{>{\bfseries}lX}
\toprule
Field & Detail \\
\midrule
Purpose & Installs newer versions of all currently installed packages \\
What it does & Downloads and installs updates without removing any existing packages \\
Expected result & Summary showing number of packages upgraded \\
Flag \texttt{-y} & Automatically answers ``yes'' to all prompts \\
\bottomrule
\end{tabularx}
\end{table}

\subsection{Command: \texttt{sudo apt full-upgrade}}

\begin{lstlisting}[style=terminal, caption={Full system upgrade}]
sudo apt full-upgrade -y
\end{lstlisting}

\begin{table}[H]
\centering
\caption{Command explanation: sudo apt full-upgrade}
\begin{tabularx}{\textwidth}{>{\bfseries}lX}
\toprule
Field & Detail \\
\midrule
Purpose & Performs a complete system upgrade including dependency changes \\
What it does & May install new packages or remove obsolete ones to satisfy dependencies \\
When to use & Recommended over \texttt{upgrade} for Kali — handles kernel and major tool updates \\
Expected result & System fully up to date; possible reboot prompt for kernel updates \\
\bottomrule
\end{tabularx}
\end{table}

\noteblock{Run the full update sequence in order: \cmd{apt update} → \cmd{apt upgrade} → \cmd{apt full-upgrade}. Skipping \texttt{update} first will upgrade to outdated cached versions.}

\subsection{Verify System Status}

\begin{lstlisting}[style=terminal, caption={Check system is up to date}]
# Check Kali version
cat /etc/os-release

# Check last update time
stat /var/lib/apt/periodic/update-success-stamp

# Check available updates (without installing)
apt list --upgradable 2>/dev/null
\end{lstlisting}

% ─────────────────────────────────────────────────────────
\section{Installing Required Tools}

\subsection{Burp Suite Community Edition}

\begin{lstlisting}[style=terminal, caption={Install Burp Suite on Kali}]
sudo apt install burpsuite -y
\end{lstlisting}

\begin{table}[H]
\centering
\caption{Burp Suite installation details}
\begin{tabularx}{\textwidth}{>{\bfseries}lX}
\toprule
Field & Detail \\
\midrule
Purpose & The primary web application security testing proxy \\
Expected result & Burp Suite appears in Applications $\to$ Web Application Analysis \\
Launch command & \texttt{burpsuite} in terminal or click application icon \\
Troubleshooting & If not found: \texttt{sudo apt install default-jdk} (Java required) \\
\bottomrule
\end{tabularx}
\end{table}

\subsection{Firefox}

\begin{lstlisting}[style=terminal, caption={Install Firefox ESR}]
sudo apt install firefox-esr -y
\end{lstlisting}

\noteblock{Firefox ESR is recommended over Chromium for Burp Suite use. It offers better proxy configuration control and FoxyProxy extension support.}

\subsection{curl}

\begin{lstlisting}[style=terminal, caption={Install and verify curl}]
sudo apt install curl -y
curl --version
\end{lstlisting}

\begin{table}[H]
\centering
\caption{curl details}
\begin{tabularx}{\textwidth}{>{\bfseries}lX}
\toprule
Field & Detail \\
\midrule
Purpose & Command-line HTTP client for quick endpoint testing \\
Expected result & \texttt{curl 8.x.x} with supported protocols listed \\
Common use & \texttt{curl -I https://target.com} to check headers \\
\bottomrule
\end{tabularx}
\end{table}

\subsection{wget}

\begin{lstlisting}[style=terminal, caption={Install wget}]
sudo apt install wget -y
\end{lstlisting}

\begin{table}[H]
\centering
\begin{tabularx}{\textwidth}{>{\bfseries}lX}
\toprule
Field & Detail \\
\midrule
Purpose & Download files from URLs — useful for fetching wordlists, tools \\
Example & \texttt{wget https://example.com/file.zip} \\
\bottomrule
\end{tabularx}
\end{table}

\subsection{git}

\begin{lstlisting}[style=terminal, caption={Install git}]
sudo apt install git -y
git --version
\end{lstlisting}

\begin{table}[H]
\centering
\begin{tabularx}{\textwidth}{>{\bfseries}lX}
\toprule
Field & Detail \\
\midrule
Purpose & Clone security tool repositories from GitHub \\
Example & \texttt{git clone https://github.com/danielmiessler/SecLists.git} \\
\bottomrule
\end{tabularx}
\end{table}

\subsection{Python 3}

\begin{lstlisting}[style=terminal, caption={Install Python 3 and pip}]
sudo apt install python3 python3-pip -y
python3 --version
pip3 --version
\end{lstlisting}

\noteblock{Python 3 is pre-installed on Kali. You mainly need pip to install security-related Python packages like \texttt{requests}, \texttt{pwntools}, \texttt{impacket}.}

\subsection{Docker}

\begin{lstlisting}[style=terminal, caption={Install Docker for running lab containers}]
sudo apt install docker.io -y
sudo systemctl enable docker
sudo systemctl start docker
sudo usermod -aG docker $USER
docker --version
\end{lstlisting}

\begin{table}[H]
\centering
\caption{Docker installation steps explained}
\begin{tabularx}{\textwidth}{>{\bfseries}lX}
\toprule
Command & Purpose \\
\midrule
\texttt{install docker.io} & Installs Docker from Kali repository \\
\texttt{systemctl enable} & Auto-start Docker on boot \\
\texttt{systemctl start} & Start Docker service now \\
\texttt{usermod -aG docker} & Allow current user to run Docker without sudo \\
\bottomrule
\end{tabularx}
\end{table}

\warnblock{Log out and back in after \texttt{usermod} for group changes to take effect. Without this, you will get \texttt{permission denied} errors running Docker commands.}

\subsection{One-Line Tool Install (All at Once)}

\begin{lstlisting}[style=terminal, caption={Install all required tools in one command}]
sudo apt update && sudo apt install -y \
  burpsuite \
  firefox-esr \
  curl \
  wget \
  git \
  python3 \
  python3-pip \
  docker.io \
  seclists \
  gobuster \
  nikto \
  nmap
\end{lstlisting}

% ─────────────────────────────────────────────────────────
\section{Essential Security Wordlists}

\begin{lstlisting}[style=terminal, caption={Install SecLists}]
sudo apt install seclists -y
ls /usr/share/seclists/
\end{lstlisting}

Key wordlist locations:

\begin{table}[H]
\centering
\caption{Important SecLists wordlist paths}
\begin{tabularx}{\textwidth}{>{\bfseries}lX}
\toprule
Path & Use Case \\
\midrule
\texttt{/usr/share/seclists/Discovery/Web-Content/} & Directory and file brute forcing \\
\texttt{/usr/share/seclists/Passwords/} & Password brute forcing \\
\texttt{/usr/share/seclists/Usernames/} & Username enumeration \\
\texttt{/usr/share/wordlists/rockyou.txt.gz} & Classic password list \\
\bottomrule
\end{tabularx}
\end{table}

\begin{lstlisting}[style=terminal, caption={Decompress rockyou.txt}]
sudo gunzip /usr/share/wordlists/rockyou.txt.gz
wc -l /usr/share/wordlists/rockyou.txt
# Should show ~14 million lines
\end{lstlisting}

\exerciseblock{
\textbf{Exercise 2.1 — System Readiness Check}\\[4pt]
Run the following commands and confirm each produces expected output:
\begin{lstlisting}[style=terminal]
burpsuite --version 2>&1 | head -1
firefox --version
curl --version | head -1
python3 --version
git --version
docker --version
\end{lstlisting}
All should return version numbers. If any fail, return to the relevant installation section above.
}
